\documentclass[11pt,a4paper]{article}
\usepackage[utf8]{inputenc}
\usepackage[T1]{fontenc}
\usepackage{amsfonts}
\usepackage[french]{babel}
\frenchbsetup{StandardLists=true} % à inclure si on utilise \usepackage[francais]{babel}
\usepackage{enumitem}
\usepackage{amssymb}
\usepackage{amsmath}
\usepackage[left=2cm,right=2cm,top=2cm,bottom=2cm]{geometry}
\usepackage{multirow}
\usepackage[dvipsnames]{xcolor}

\usepackage[T1]{fontenc}
\usepackage{fouriernc}
\usepackage{graphicx}

\usepackage{setspace}
\singlespacing

\usepackage{multicol}

\usepackage{fancyhdr}
\pagestyle{fancy}
\renewcommand\headrulewidth{1pt}
\fancyhead[L]{Lycée Denis Diderot }
\fancyhead[R]{Classe 2BTS}
\setcounter{page}{1}\fancyfoot[C]{Contrôle 4 - Page \thepage}
\fancyfoot[R]{\today}
\fancyfoot[L]{Alexis RUIZ}
\usepackage{multirow,tabularx,booktabs}
\newcommand{\cadreblanc}[1]{%
	\medskip\framebox[\linewidth][c]{%
		\rule{0mm}{#1mm}}\par
	\medskip}

\renewcommand{\arraystretch}{2}
\newcommand*\acocher{\quitvmode{\fboxsep0pt \fboxrule0.8pt \fbox{\vrule height1.8ex depth.1ex width0pt \vrule height0pt depth0pt width1.9ex }}}
\DeclareMathOperator{\e}{e}

\setcounter{Exo}{0}

\begin{document}
	
	
	\setlength{\parindent}{0cm}
	
	
	\bigskip
	
	\center \LARGE{Devoir Surveillé  - Équations différentielles (2 h)}\
	\\[0,3cm]
	
	\normalsize
	
	\hrule\
	\\[0,2cm]
	
	\flushleft \textbf{NOM :\\
		Prénom:\\[0,5cm]
	}
	\hrule\
	\\[0,2cm]
	
	%%%%%%%%%%%%%%%%%%%%%%%%%%%%%%%%%%%%%%%%%%%%%%%%%%%%%%%%%%
	
	
	\fcolorbox[gray]{0}{0.9}{\textbf{Exercice 1}} Extrait sujet Groupement B - 2012\\[2mm]
	On considère l'équation différentielle (E) :\[
	y' + 2y = -5\text{e}^{-2x}\]
	où $y$ est une fonction de la variable réelle $x$, définie
	et dérivable sur $\mathbb{R}$ et $y'$ sa fonction dérivée.
	\\
	\begin{enumerate}
		\item Résoudre sur $\mathbb{R}$ l'équation différentielle (E$_{0}$):
		\[ y' + 2y = 0\]
		
		\vspace{-5mm}
		\cadreblanc{25}
		
		
		\item Soit $g$ la fonction définie sur $\mathbb{R}$ par 
		\[g(x) = -5x\text{e}^{-2x}\]
		Démontrer que $g$ est une solution particulière de l'équation différentielle (E).
		
		\cadreblanc{50}
		
		\item En déduire l'ensemble des solutions de l'équation différentielle (E).
		
		\cadreblanc{20}
		
		\item Déterminer la solution particulière $f$ de l'équation (E) qui vérifie la condition \[f(0)=1\]
		
		\vspace{-5mm}
		
		\cadreblanc{20}
		
	\end{enumerate} 
	
	
	
	
	\vfill 
	
	\newpage
	
	%%%%%%%%%%%%%%%%%%%%%%%%%%%%%%%%%%%%%%%%%%%%%%%%%%%%%%%%%%%
	
	
	\fcolorbox[gray]{0}{0.9}{\textbf{Exercice 2}}  On considère l'équation différentielle (E) :\[
	y' - 2y = -x^{2}\]
	où $y$ est une fonction de la variable réelle $x$, définie
	et dérivable sur $\mathbb{R}$ et $y'$ sa fonction dérivée.
	\\
	\begin{enumerate}
		\item Résoudre sur $\mathbb{R}$ l'équation différentielle (E$_{0}$):
		\[ y' - 2y = 0\]
		\vspace{-5mm}
		\cadreblanc{25}
		
		\item Déterminer une solution particulière de l’équation différentielle (E) sous forme d’un polynôme. $$y_p(x)=Ax^2+Bx+C$$
		
		\cadreblanc{80}
		
		\begin{multicols}{2}
			\item En déduire l'ensemble des solutions de l'équation différentielle (E).
			
			\cadreblanc{45}
			\item Déterminer la solution particulière $f$ de l'équation (E) qui vérifie la condition \[f(0)=1\]
			\vspace{-11}
			\cadreblanc{35}
		\end{multicols}
	\end{enumerate} 
	
	\newpage
	
	\fcolorbox[gray]{0}{0.9}{\textbf{Exercice 3}} On considère l'équation différentielle (E) :\[
	y' -  y = (3x-1) \times \text{e}^{x}\]
	où $y$ est une fonction de la variable réelle $x$, définie
	et dérivable sur $\mathbb{R}$ et $y'$ sa fonction dérivée.
	\\
	\begin{enumerate}
		\item Résoudre sur $\mathbb{R}$ l'équation différentielle (E$_{0}$):
		\[ y' - y = 0\]
		\vspace{-5mm}
		\cadreblanc{25}
		
		\item Déterminer une solution particulière de l’équation différentielle (E) sous la forme $$y_p(x)= (Ax^2+Bx+C) \e^x$$.
		
		\cadreblanc{80}
		
		\begin{multicols}{2}
			\item En déduire l'ensemble des solutions de l'équation différentielle (E).
			
			\cadreblanc{45}
			\item Déterminer la solution particulière $f$ de l'équation (E) qui vérifie la condition \[f(0)=-1\]
			\vspace{-11}
			\cadreblanc{35}
		\end{multicols}
	\end{enumerate} 
	
	
	
\end{document}